\documentclass[10pt,a4paper]{amsart}
\usepackage[margin=19mm]{geometry}
\usepackage{amsmath,amssymb,mathtools}
\usepackage[scr=rsfs,cal=euler]{mathalfa}
\usepackage{xfrac}
\usepackage[unicode,hidelinks,bookmarks=false]{hyperref}
\newcommand{\ie}{i.e.\ }
\newcommand{\cf}{cf.\ }
\newcommand{\resp}{respectively\ }
\newcommand{\Subsection}{\subsection}
\providecommand{\nobreakdash}{\nobreak\hbox{-}}
\setlength{\emergencystretch}{2em}
\multlinegap0pt
\usepackage{bm}
\usepackage{bbm}
\usepackage{dsfont}
\usepackage{xspace}
\usepackage{cancel}
\usepackage{mathtools}
\usepackage{amsthm}
\makeatletter
\@ifundefined{definition}{\newtheorem{definition}{Definition}}{}
\@ifundefined{theorem}{\newtheorem{theorem}[definition]{Theorem}}{}
\makeatother
\title[Generalized Discrete Orlicz-Morrey Spaces]{Generalized Discrete Orlicz-Morrey Spaces}
\author{Hukmashabiyya Ariq Gumilar, Al Azhary Masta, Siti Fatimah}
\date{Source: arXiv:2511.04004v1 (CC BY 4.0)}
\begin{document}
\maketitle

\noindent\emph{Source and licence.} Generalized Discrete Orlicz-Morrey Spaces. Version arXiv:2511.04004v1 (CC BY 4.0), distributed under the Creative Commons Attribution 4.0 International licence, as recorded in the arXiv licence metadata for this exact version and shown on the abstract page https://arxiv.org/abs/2511.04004v1. Full rights notice: licenses/arxiv-2511.04004-CC-BY-4.0.txt.\par\smallskip

\noindent\textit{Editorial scope.} Counted unit: the Theorem whose source label is \texttt{None} in the article (source0.tex lines 397--399, source label \texttt{None}), delivered together with the article's own complete original proof of it (source0.tex lines 401--525, one \texttt{proof} environment of 125 lines). No mathematics is written, repaired or re-derived: statement and proof are the article's own text line for line. Required context retained uncounted: none. Every definition, display and label of those lines is retained unchanged. Omitted from this package and not redistributed: 1--396, 400--400, 526--587. Nothing inside a retained excerpt is omitted or altered: every retained body is delivered exactly as the article's lines are written, so no omission receipt of any kind accompanies this package. Citations: the retained excerpts carry no citation command at all, so no bibliography entry of the article is transcribed and no reference is redistributed. Reference adaptation: none was needed and none was made; every reference inside the delivered unit resolves inside it, so no unresolved reference or citation is produced. Printed slips kept literal: no printed word, symbol, hypothesis or number of the counted statement or of its proof was altered. Layout and class: the article's own class is \texttt{article} with packages including inputenc, fontenc, lmodern, babel, amsmath, amssymb, amsthm, graphicx, hyperref, geometry, cite, setspace, enumerate, caption; the wrapper is the lane's pinned \texttt{amsart} editorial wrapper at 10pt on A4, which restates the theorem-like environments the retained text needs and adds the article's own macro definitions that the retained text needs (none), with extra wrapper packages (bm, bbm, dsfont, xspace, cancel, mathtools). The delivered numbering is the unit's own. Assets: no retained span contains a figure environment, a graphics-inclusion command, a PDF-page inclusion, a table or a TikZ picture; no image file is copied into this package, the package declares no asset dependency and no image file is redistributed with it. Licence: version arXiv:2511.04004v1 of this article is distributed under the Creative Commons Attribution 4.0 International licence, as recorded in the arXiv licence metadata for this exact version and shown on the abstract page https://arxiv.org/abs/2511.04004v1. A full rights notice is delivered at licenses/arxiv-2511.04004-CC-BY-4.0.txt. Characters: every retained excerpt is pure ASCII, so the delivered engine, XeLaTeX with its default Latin Modern text font, reports no missing character.
\begin{theorem}
    The nonnegative function $\left\lVert \cdot \right\rVert _{\ell_{\phi,\Phi_s}}$ is quasi-norm.
\end{theorem}

\begin{proof}
    Consider any sequences $x=\left(x_k\right)$ and $y=\left(y_k\right)$ where $x, y \in \ell_{\phi, \Phi_s}$. To demonstrate that $\left\lVert x \right\rVert _{\ell_{\phi,\Phi_s}}$ is a quasi-norm function, it is necessary to verify the following four conditions based on Definition 2.15:

    \begin{enumerate}
        \item It will be shown that $\left\lVert x \right\rVert _{\ell_{\phi,\Phi_s}}\geq 0$.
        Let $A = \left\{b>0 \;\middle|\; \frac{\phi\left(2N+1\right)}{\left|S_{m,N}\right|}\sum_{k\in S_{m,N}}\Phi_s\left(\frac{\left|x_k\right|}{b}\right) \leq 1\right\}$. Since $\left\lVert x \right\rVert _{\phi, \Phi_s, m, N}=\inf A$, by the definition of the infimum, $0 \leq \|x\|_{\phi, \Phi_s, m, N} \leq b$ for every $b \in A$. Thus, $\left\lVert x \right\rVert _{\phi, \Phi_s, m, N}\geq 0$. Therefore, by the definition of $\left\lVert \cdot \right\rVert _{\ell_{\phi,\Phi_s}}$,
        \begin{eqnarray*}
            \left\lVert x \right\rVert _{\ell_{\phi,\Phi_s}}=\sup_{m \in \mathbb{Z}, N \in \mathbb{N}_0} \left\lVert x \right\rVert _{\phi, \Phi_s, m, N}\geq 0
        \end{eqnarray*}
        Hence, it is proved that $\left\lVert x \right\rVert _{\ell_{\phi,\Phi_s}}\geq 0$.
        
        \item It will be shown that $\left\lVert x \right\rVert _{\ell_{\phi,\Phi_s}}=0$ if and only if $x=0$.
        \begin{enumerate}
            \item [$(\Leftarrow)$] Assume $x = \left(x_k\right) = 0$. It will be demonstrated that $\left\lVert x \right\rVert _{\ell_{\phi,\Phi_s}}=0$. Take any $b>0$, $m \in \mathbb{Z}$, and $N \in \mathbb{N}_0$, observe that 
            \begin{eqnarray*}
                \frac{\phi\left(2N+1\right)}{\left|S_{m,N}\right|}\sum_{k\in S_{m,N}}\Phi_s\left(\frac{\left|x_k\right|}{b}\right)=0 \leq 1
            \end{eqnarray*}
            Therefore, $b \in A$, and it follows that $\left\lVert x \right\rVert _{\phi, \Phi_s, m, N}\leq b$ for every $b>0$, implying $\left\lVert x \right\rVert _{\phi, \Phi_s, m, N}=0$. By definition of $\left\lVert x \right\rVert _{\ell_{\phi,\Phi_s}}$,
            \begin{eqnarray*}
                \left\lVert x \right\rVert _{\ell_{\phi,\Phi_s}}&=&\sup_{m \in \mathbb{Z}, N \in \mathbb{N}_0} \left\lVert x \right\rVert _{\phi, \Phi_s, m, N}\\
                   &=&0.
            \end{eqnarray*}
               
            \item[($\Rightarrow$)] Assume $\left\lVert x \right\rVert _{\ell_{\phi,\Phi_s}}=0$. Observe that $\left\lVert x\right\rVert _{\ell_{\phi,\Phi_s}} =\sup_{m\in\mathbb{Z}, N\in\mathbb{N}_0} \left\lVert x\right\rVert _{\phi,\Phi_s, m, N}=0$. It means $\left\lVert x\right\rVert _{\phi,\Phi_s, m, N}=0$. By Lemma 3.6, it can be obtained that $\frac{\phi\left(2N+1\right)}{\left|S_{m,N}\right|}\sum_{k\in S_{m,N}}\Phi_s\left(a\left|x_k\right|\right)=0$. Therefore, it must be that $\Phi_s\left(a\left|x_k\right|\right) = 0$. Consequently, $x_k = 0$ for all $k \in S_{m,N}$, which implies $x = \left(x_k\right)_{k\in \mathbb{Z}} = 0$. 
            It is proven that $\left\lVert x \right\rVert _{\ell_{\phi,\Phi_s}}=0$ if and only if $x=0$..
        \end{enumerate}
        
        \item It will be shown that $\left\lVert ax \right\rVert _{\ell_{\phi,\Phi_s}}=\left|a\right|\cdot\left\lVert x \right\rVert _{\ell_{\phi,\Phi_s}}$ for every $a \in \mathbb{R}$.
        
        Case 1: $a = 0$.\\
        Observe that
        \begin{eqnarray*}
            \left\lVert ax \right\rVert _{\ell_{\phi,\Phi_s}}&=&\left\lVert 0\cdot x \right\rVert _{\ell_{\phi,\Phi_s}}\\
            &=&\left\lVert 0\cdot x \right\rVert _{\ell_{\phi,\Phi_s}}\\
            &=&0\\
            &=&\left|0\right|\cdot\left\lVert x \right\rVert _{\ell_{\phi,\Phi_s}}.
        \end{eqnarray*}
        case 2: $a \neq 0$.\\
        Take any $a \in \mathbb{R} \setminus \left\{0\right\}$. Note that
        \begin{eqnarray*}
            \left\lVert ax \right\rVert _{\phi, \Phi_s, m, N}&=&\inf \left\{b > 0 \;\middle|\; \frac{\phi\left(2N+1\right)}{\left|S_{m,N}\right|}\sum_{k\in S_{m,N}}\Phi_s\left(\frac{\left|a\cdot x_k\right|}{b}\right) \leq 1\right\}\\
            &=&\inf \left\{b > 0 \;\middle|\; \frac{\phi\left(2N+1\right)}{\left|S_{m,N}\right|}\sum_{k\in S_{m,N}}\Phi_s\left(\frac{\left|a\right|\cdot\left|x_k\right|}{b}\right) \leq 1\right\}\\
            &=&\inf \left\{b > 0 \;\middle|\; \frac{\phi\left(2N+1\right)}{\left|S_{m,N}\right|}\sum_{k\in S_{m,N}}\Phi_s\left(\frac{\left|x_k\right|}{\frac{b}{\left|a\right|}}\right) \leq 1\right\}\\
        \end{eqnarray*}
        Using the substitution $c = \frac{b}{\left|a\right|}$, it follows that
        \begin{eqnarray*}
            \left\lVert ax \right\rVert _{\phi, \Phi_s, m, N}&=&\inf \left\{c\left|a\right|> 0 \;\middle|\; \frac{\phi\left(2N+1\right)}{\left|S_{m,N}\right|}\sum_{k\in S_{m,N}}\Phi_s\left(\frac{\left|x_k\right|}{c}\right) \leq 1\right\}\\
            &=&\left|a\right|\inf \left\{c > 0 \;\middle|\; \frac{\phi\left(2N+1\right)}{\left|S_{m,N}\right|}\sum_{k\in S_{m,N}}\Phi_s\left(\frac{\left|x_k\right|}{c}\right) \leq 1\right\}\\
            &=&\left|a\right|\cdot\left\lVert x \right\rVert _{\phi, \Phi_s, m, N}.
        \end{eqnarray*}
        Thus,
        \begin{eqnarray*}
            \left\lVert ax \right\rVert _{\ell_{\phi,\Phi_s}}&=&\sup_{m \in \mathbb{Z}, N \in \mathbb{N}_0} \left\lVert a x\right\rVert_{\phi, \Phi_s, m, N}\\
            &=& \left|a\right|\cdot\sup_{m \in \mathbb{Z}, N \in \mathbb{N}_0} \left\lVert x\right\lVert_{\phi, \Phi_s, m, N}\\ 
            &=& \left| a\right| \left\lVert x\right\lVert_{\ell_{\phi, \Phi_s}}.
        \end{eqnarray*}
        Therefore, it is proven that $\left\lVert ax \right\rVert _{\ell_{\phi,\Phi_s}}=\left| a\right| \left\lVert x\right\lVert_{\ell_{\phi, \Phi_s}}$. for every $a\in \mathbb{R}$.
        
        \item It will be shown that there exists $C\geq 1$ such that
        \begin{eqnarray*}
            \left\lVert x+y \right\rVert _{\ell_{\phi,\Phi_s}}\leq C^{\frac{1}{s}} \left(\left\lVert x \right\rVert _{\ell_{\phi, \Phi_s}} + \left\lVert y \right\rVert _{\ell_{\phi, \Phi_s}} \right). \forall s\in \left(0,1\right].
        \end{eqnarray*}
        Case 1: $x=0$ and $y=0$.\\
        Since $x=0$ and $y=0$, it follows that $\left\lVert x \right\rVert _{\ell_{\phi,\Phi_s}}=0$ and $\left\lVert y \right\rVert _{\ell_{\phi,\Phi_s}}=0$.
        Therefore,
        \begin{eqnarray*}
            \left\lVert x+y \right\rVert _{\ell_{\phi,\Phi_s}} &=&\left\lVert 0 \right\rVert _{\ell_{\phi,\Phi_s}}\\
            &=& 0\\
            &=& C^{\frac{1}{s}}\left(0+0\right)\\
            &=& C^{\frac{1}{s}}\left(\left\lVert x \right\rVert _{\ell_{\phi,\Phi_s}}+\left\lVert y \right\rVert _{\ell_{\phi,\Phi_s}}\right). 
        \end{eqnarray*}
        Case 2: $x\neq0$ and $y=0$ or vise versa.\\  
        Without loss of generality, assume $x\neq0$ and $y=0$. Choose $C=1$, so 
        \begin{eqnarray*}
            \left\lVert x+y \right\rVert _{\ell_{\phi,\Phi_s}} &=&\left\lVert x+0 \right\rVert _{\ell_{\phi,\Phi_s}}\\
            &=& 1\left(\left\lVert x \right\rVert _{\ell_{\phi,\Phi_s}}+0\right)\\ &=&C^{\frac{1}{s}}\left(\left\lVert x \right\rVert _{\ell_{\phi,\Phi_s}}+\left\lVert y \right\rVert _{\ell_{\phi,\Phi_s}}\right).
        \end{eqnarray*}
        Case 3: $x\neq 0$, $y\neq 0$.\\
        Let $B = \left\{b > 0 \;\middle|\; \frac{\phi\left(2N+1\right)}{\left|S_{m,N}\right|}\sum_{k\in S_{m,N}} \Phi_s \left(\frac{\left|x_k + y_k\right|}{b}\right) \leq 1\right\}$, $X=\left\lVert x \right\rVert _{\phi, \Phi_s, m, N}$, and $Y=\left\lVert y\right\rVert _{\phi, \Phi_s, m, N}$.
        Define $C = \frac{X^s + Y^s}{(X + Y)^s}$. Observe that
        \begin{eqnarray*}
            \frac{\phi\left(2N+1\right)}{\left|S_{m,N}\right|}\sum_{k\in S_{m,N}}\Phi_s\left(\frac{\left|x_k+y_k\right|}{C^{\frac{1}{s}} \left(X+Y\right)}\right) &\leq  \frac{\phi\left(2N+1\right)}{\left|S_{m,N}\right|}\sum_{k\in S_{m,N}}\Phi_s\left(\frac{\left|x_k\right|+\left|y_k\right|}{C^{\frac{1}{s}} \left(X+Y\right)}\right)\\
            &=  \frac{\phi\left(2N+1\right)}{\left|S_{m,N}\right|}\sum_{k\in S_{m,N}}\Phi_s\left(\frac{\left|x_k\right|}{C^{\frac{1}{s}} \left(X+Y\right)}\frac{X}{X} \right.\\
            & \left.+\frac{\left|y_k\right|}{C^{\frac{1}{s}} \left(X+Y\right)}\frac{Y}{Y}\right). 
        \end{eqnarray*}
        Because of $\left(\frac{X}{C^{\frac{1}{s}} \left(X+Y\right)}\right)^s+\left(\frac{Y}{C^{\frac{1}{s}} \left(X+Y\right)}\right)^s=1$, by \emph{s}-convex definition,
        \begin{multline*}
            \frac{\phi(2N+1)}{|S_{m,N}|} \sum_{k\in S_{m,N}} 
            \Phi_s\!\left(
               \frac{|x_k|}{C^{1/s}(X+Y)}\frac{X}{X}
               + \frac{|y_k|}{C^{1/s}(X+Y)}\frac{Y}{Y}
            \right)
            \\
            \leq
            \left(\frac{X}{C^{1/s}(X+Y)}\right)^s 
               \frac{\phi(2N+1)}{|S_{m,N}|}\sum_{k\in S_{m,N}} \Phi_s\!\left(\frac{|x_k|}{X}\right)
            \\
            + \left(\frac{Y}{C^{1/s}(X+Y)}\right)^s 
               \frac{\phi(2N+1)}{|S_{m,N}|}\sum_{k\in S_{m,N}} \Phi_s\!\left(\frac{|y_k|}{Y}\right)
            \\
            \leq
            \left(\frac{X}{C^{1/s}(X+Y)}\right)^s
            + \left(\frac{Y}{C^{1/s}(X+Y)}\right)^s
            = 1.
        \end{multline*}


        It means $C^{\frac{1}{s}}\left(X+Y\right)=C^{\frac{1}{s}}\left(\left\lVert x \right\rVert _{\phi, \Phi_s, m, N}+\left\lVert y \right\rVert _{\phi, \Phi_s, m, N}\right)\in B$. Since $\left\lVert x+y \right\rVert _{\phi, \Phi_s, m, N}$ is infimum of $B$, then $\left\lVert x+y \right\rVert _{\phi, \Phi_s, m, N}\leq C^{\frac{1}{s}}\left(X+Y\right)=C^{\frac{1}{s}}\left(\left\lVert x \right\rVert _{\phi, \Phi_s, m, N}+\left\lVert y \right\rVert _{\phi, \Phi_s, m, N}\right)$.
        By taking the supremum over $m \in \mathbb{Z}$ and $N \in \mathbb{N}_0$, it can be obtained that $\left\lVert x+y \right\rVert _{\ell_{\phi,\Phi_s}}\leq C^{\frac{1}{s}} \left(\left\lVert x \right\rVert _{\ell_{\phi, \Phi_s}} + \left\lVert y \right\rVert _{\ell_{\phi, \Phi_s}} \right)$.
        Furthermore, it will be demonstrated that $1 \leq C < 2$. For every $s \in \left(0,1\right]$, observe that
        \begin{eqnarray*}
            \frac{X}{X + Y} \leq \left(\frac{X}{X + Y}\right)^s < 1
        \end{eqnarray*}
        and
        \begin{eqnarray*}
            \frac{Y}{X + Y} \leq \left(\frac{Y}{X + Y}\right)^s < 1.
        \end{eqnarray*}
        As a result,
        \begin{eqnarray*}
            1 = \frac{X + Y}{X + Y} \leq \frac{X^s + Y^s}{(X + Y)^s} = C < 1 + 1 = 2.
        \end{eqnarray*}
        Thus, it has been established that \(1 \leq C < 2\), ensuring that $\left\lVert x+y \right\rVert _{\ell_{\phi,\Phi_s}}\leq C^{\frac{1}{s}} \left(\left\lVert x \right\rVert _{\ell_{\phi, \Phi_s}} + \left\lVert y \right\rVert _{\ell_{\phi, \Phi_s}} \right). \forall s\in \left(0,1\right]$.  
    \end{enumerate} 
    Since conditions (1) through (4) are satisfied, it is established that $\left\lVert x \right\rVert _{\ell_{\phi, \Phi_s}}$ is a quasi-norm on the Orlicz-Morrey-s sequence space $\ell_{\phi, \Phi_s}$.
\end{proof}

\end{document}
